\chapter{Appendix 2: An example of User-level Security Policy}


Ultimately, computer security is the user's responsibility. Each user
must be alert to possible breaches in security and adhere to the
security regulations that have been established within their company.
The security practices listed here are not inclusive, but rather
designed to remind them and raise their awareness towards securing
their information resources:

\begin{htmllist}
  \htmlitemmark{RedBall}

\item[Protect Your Equipment]

Keep it in a secure environment

Keep food, drink, and cigarettes AWAY from it

Know where the fire suppression equipment is located and know how to
use it


\item[Protect Your Area]

Keep unauthorized people AWAY from your equipment and data

Challenge strangers in your area


\item[Protect Your Password]

Never write it down or give it to anyone

Don't use names, numbers or dates which are personally identified with
you

Change it often, but change it immediately if you think it has been
compromised


\item[Protect Your Files]

Don't allow unauthorized access to your files and data

NEVER leave your equipment unattended with your password activated -
SIGN OFF!


\item[Protect Against Viruses]

Don't use unauthorized software

Back up your files before implementing ANY new software


\item[Lock Up Storage Media Containing Sensitive Data]

If the data or information is sensitive or critical to your operation,
lock it up!


\item[Back Up Your Data]

Keep duplicates of your sensitive data in a safe place, out of your
immediate area

Back it up as often as necessary


\item[Report Security Violations]

Tell your manager if you see any unauthorized changes to your data

Immediately report any loss of data or programs, whether automated or
hard copy

\end{htmllist}
